Figure~\ref{fig:architecture} summarizes the architecture of the \omniBase{} models.
We extend Jina Embeddings v5 Text from text-only embedding to vision and audio by adding scale-matched Qwen3.5 vision encoders\footnote{\omniSmall{} uses \texttt{Qwen/Qwen3.5-2B}; \omniNano{} uses \texttt{Qwen/Qwen3.5-0.8B}.} and the Qwen2.5-Omni audio encoder to the same text-sequence backbone.
We chose encoders from trained multimodal language systems rather than bare perceptual encoders such as SigLIP2 or Whisper-large because prior work shows that visual and audio features need explicit language-space alignment or natural-language supervision before they transfer reliably to text-conditioned multimodal tasks~\cite{comp,clap,qwen3.5,qwen2.5-omni}.
The text processing path of \omniBase{} shares the frozen text-encoder weights of Jina Embeddings v5 Text bit-for-bit: token embeddings pass through the frozen text transformer, the inherited task LoRA adapter is applied, and the final embedding is produced by last-token pooling and L2 normalization.

\subsection{Projectors}
\label{sec:method:projectors}

\omniBase{} uses image and audio encoders extracted from Qwen3.5 and Qwen2.5-Omni, respectively. Because their output dimensions do not match Jina Embeddings v5 Text's input, we replace the source projection layers with new projectors that map into the text hidden space.
For audio, we inserted a randomly-initialized \texttt{fc\_audio} layer that projects the encoder's native $1280$ dimension output into \omniSmall{}'s $1024$-dimension input space and \omniNano{}'s $768$-dimension one.

We write each fully connected layer as the same affine map
\begin{equation*}
\ell_{W,\mathbf{b}}(\mathbf{x}) = W\mathbf{x} + \mathbf{b},
\label{eq:linear_layer}
\end{equation*}
with layer-specific weights and bias.
Thus \texttt{fc\_vision\_1} is $\ell_{W_{\text{v1}},\mathbf{b}_{\text{v1}}}$, \texttt{fc\_vision\_2} is $\ell_{W_{\text{v2}},\mathbf{b}_{\text{v2}}}$, and \texttt{fc\_audio} is $\ell_{W_{\text{aud}},\mathbf{b}_{\text{aud}}}$.

For vision, the Qwen3.5 visual projector converts ViT patch tokens into text-token features by applying LayerNorm, a $2{\times}2$ spatial merge, \texttt{fc\_vision\_1}, GELU, and \texttt{fc\_vision\_2}.
Here, LayerNorm denotes feature normalization on the ViT patch tokens.
The $2{\times}2$ spatial merge is a fixed space-to-depth (pixel-unshuffle) rearrangement that concatenates four neighboring patch embeddings into one $4d_{\text{vit}}$ vector, reducing the spatial token count by $4\times$; it is the inverse direction of pixel shuffle/sub-pixel rearrangement~\cite{pixel-shuffle} and follows Qwen's visual-merger design~\cite{qwen2-vl,qwen3.5}.
For each group of four neighboring patch tokens $\mathbf{V}_i = [\mathbf{v}_{i,1}, \ldots, \mathbf{v}_{i,4}] \in \mathbb{R}^{4 \times d_{\text{vit}}}$, the vision projector produces
\begin{equation*}
\begin{aligned}
\mathbf{m}^{(i)}_{\text{vis}} &=
\bigl[
\text{LayerNorm}(\mathbf{v}_{i,1});
\ldots;
\text{LayerNorm}(\mathbf{v}_{i,4})
\bigr] \in \mathbb{R}^{4d_{\text{vit}}},\\
\mathbf{z}^{(i)}_{\text{vis}} &=
\text{GELU}\!\left(\ell_{W_{\text{v1}},\mathbf{b}_{\text{v1}}}(\mathbf{m}^{(i)}_{\text{vis}})\right),\\
\mathbf{h}^{(i)}_{\text{vis}} &=
\ell_{W_{\text{v2}},\mathbf{b}_{\text{v2}}}(\mathbf{z}^{(i)}_{\text{vis}}),
\qquad i = 1,\ldots,N_{\text{vis}}.
\end{aligned}
\label{eq:vis_proj}
\end{equation*}

Only \texttt{fc\_vision\_2} performs the dimension-specific projection into a text hidden space: in the 2B source checkpoint it maps $4096{\to}2048$ into the Qwen3.5-2B text hidden dimension, and in the 0.8B source checkpoint it maps $3072{\to}1024$ into the Qwen3.5-0.8B text hidden dimension.
These targets do not match Small's $1024$-dimensional or Nano's $768$-dimensional Jina text backbone, so we keep LayerNorm and \texttt{fc\_vision\_1} frozen but replace \texttt{fc\_vision\_2} with a randomly initialized $4096{\to}1024$ layer for Small and $3072{\to}768$ layer for Nano.

Let $\mathbf{A} = [\mathbf{a}_1, \ldots, \mathbf{a}_K] \in \mathbb{R}^{K \times 1280}$ denote the frozen Qwen2.5-Omni audio encoder states for an input with $K$ audio tokens.
Each audio token is independently projected into the Jina text hidden dimension by \texttt{fc\_audio}
\begin{equation*}
\mathbf{h}^{(i)}_{\text{aud}} =
\ell_{W_{\text{aud}},\mathbf{b}_{\text{aud}}}(\mathbf{a}_i),
\qquad i = 1,\ldots,K,
\label{eq:aud_proj}
\end{equation*}
where $W_{\text{aud}} \in \mathbb{R}^{d_{\text{text}} \times 1280}$ and $d_{\text{text}} \in \{1024,768\}$ for Small and Nano.

\subsection{Input Sequence Construction}
\label{sec:method:input-sequence}

Each input is serialized as one sequence of tokens.
Text remains ordinary text tokens; non-text modalities are represented by placeholder runs inside modality delimiters.
An image is encoded as
\begin{equation*}
\texttt{<|vision\_start|>} \;\; \underbrace{\texttt{<|image\_pad|>} \times N}_{\text{visual slots}} \;\; \texttt{<|vision\_end|>}
\label{eq:img_seq}
\end{equation*}
with $N$ visual slots.
An audio input is encoded as
\begin{equation*}
\texttt{<|audio\_start|>} \;\; \underbrace{\texttt{<|audio\_pad|>} \times K}_{\text{audio slots}} \;\; \texttt{<|audio\_end|>}
\label{eq:aud_seq}
\end{equation*}
with $K$ audio slots.
A video is a concatenation of one visual segment per sampled frame:
\begin{equation*}
\big\Vert_{f=1}^{F}
\left(
\texttt{<|vision\_start|>} \;\; \underbrace{\texttt{<|video\_pad|>} \times S_f}_{\text{frame } f \text{ slots}} \;\; \texttt{<|vision\_end|>}
\right),
\label{eq:video_seq}
\end{equation*}
where $\Vert$ denotes sequence concatenation.
If a video contains an audio track, the extracted audio segment precedes the frame sequence:
\begin{equation*}
\mathbf{s}_{\text{aud}} \Vert \mathbf{s}_{\text{vid}}.
\label{eq:video_audio_seq}
\end{equation*}
Here, $\mathbf{s}_{\text{aud}}$ is the audio sequence above and $\mathbf{s}_{\text{vid}}$ is the video-frame sequence.
For mixed-modality inputs, text spans and modality segments are concatenated in document order.

\subsection{Trainable Parameters}
\label{sec:method:trainable-parameters}

The trainable set is \texttt{fc\_vision\_2}, \texttt{fc\_audio}, and the modality-delimiter embeddings.
\omniSmall{} learns the vision and audio start/end delimiter embeddings used in Section~\ref{sec:method:input-sequence}; \omniNano{} learns only the audio start/end delimiter embeddings.
The image, video, and audio placeholder positions are overwritten by projected encoder features rather than learned as standalone token embeddings.
Projector and delimiter-token training is run separately for retrieval, text-matching, clustering, and classification, while the text transformer, encoder towers, LayerNorm/\texttt{fc\_vision\_1} vision-projector weights, and inherited LoRA adapters stay frozen.
The base package stores four such task-specific sets alongside the inherited LoRA adapters.

\subsection{Dynamic Weight Loading}
\label{sec:method:dynamic-loading}

Jina Embeddings v5 Text already uses dynamic adapter selection to route retrieval, classification, clustering, and text-matching inputs through the corresponding task adapter.
We extend the same task-selection mechanism to the multimodal weights: the selected task variant determines which LoRA adapter, \texttt{fc\_vision\_2}, \texttt{fc\_audio}, and learned special text-token embeddings are loaded or activated.
The task-specific projector and delimiter-token weights therefore follow the same task-specific variation as Jina Embeddings v5 Text.
Separately, the model exposes a modality attribute that controls which frozen modality towers are instantiated: text-only loading omits both vision and audio towers, vision-only loading omits the audio tower and \texttt{fc\_audio}, audio-only loading omits the vision tower and vision projector, and omni loading keeps both vision and audio towers.
